\subsection{Flat anticanonical curves}\label{subsec:nef-flat-anticanonical-curves}

A flat neighborhood of a divisor gives a closed semipositive form whose curvature is confined to a thin annulus around the divisor. On a surface this form has rank at most one. The maximum-rank criterion then applies to its positive variable multiples, even when the divisor line has no moving positive multiple. We first construct the flat neighborhood for a square-zero anticanonical curve.

\begin{lemma}\label{lem:flat-elliptic-derivative-inverse}
Let $D$ be an elliptic curve with a flat Hermitian metric, let $\vartheta$ be a parallel unit $(1,0)$-form, and let $v$ be its dual vector field. For every unitary flat line bundle $L$ on $D$ and every smooth $L$-valued $(0,1)$-form $u$, there is a unique smooth section $f$ orthogonal to the parallel sections such that
\[
 \bar\partial_L f=\iota_v\partial_Lu.
\]
If $\Pi_L^{0,1}$ is the orthogonal projection onto the parallel $L$-valued $(0,1)$-forms, then, for every flat Sobolev norm $H^k$,
\[
 \|f\|_{H^k}=\|(I-\Pi_L^{0,1})u\|_{H^k}
 \leq\|u\|_{H^k}.
\]
The bound is independent of the holonomy of $L$.
\end{lemma}
\begin{proof}
Put $P=\iota_v\partial_L$. Flatness gives $P=\nabla_v$, $P^*=-\nabla_{\bar v}$, and $\Delta_L=P^*P=PP^*$ on functions and $(0,1)$-forms. Let $G_q$ be the inverse of $\Delta_L$ on the orthogonal complement of its kernel in degree $q$, extended by zero on that kernel. The operators $P$ and $P^*$ commute with $G_q$. Set
\[
 f=\bar\partial_L^*G_1Pu.
\]
The form $Pu$ is orthogonal to the parallel forms, so $\bar\partial_Lf=Pu$; also $f$ is orthogonal to the parallel sections. Uniqueness follows from the kernel of $\bar\partial_L$. Moreover,
\[
 \|f\|_{L^2}^2
 =\langle Pu,G_1Pu\rangle
 =\langle u,(I-\Pi_L^{0,1})u\rangle.
\]
Covariant differentiation along the global parallel real frame commutes with these operators. Applying the identity to every derivative proves the $H^k$ assertion. In particular, this argument never divides by a lower bound for a nonzero eigenvalue; the derivative in $Pu$ cancels that small eigenvalue.
\end{proof}

For the neighborhood argument, let $S$ be a complex surface, possibly noncompact, and let $D\subset S$ be a smooth compact connected anticanonical divisor with $D^2=0$. Adjunction makes $D$ elliptic. Its normal line $N$ has degree zero, hence admits a unitary flat metric. Let $\pi:N\to D$ be its projection, $N^\times$ its punctured total space, and $J_N$ its complex structure. The reciprocal anticanonical section gives a meromorphic two-form $\Omega_S$, nowhere zero on $S\setminus D$, with simple pole along $D$. Denote its nonzero residue by $\vartheta_D$.

Choose parallel unitary local frames of $N$ with fiber coordinate $w$. On $N^\times$ put
\[
 \xi=w\partial_w,\qquad \Theta=\frac{\mathrm dw}{w},\qquad
 \vartheta=\pi^*\vartheta_D,\qquad
 z=-\log|w|^2,\qquad \Omega_0=\Theta\wedge\vartheta.
\]
These objects are global. Give $D$ the flat metric making $\vartheta_D$ unit, and let $v$ be the parallel vector field dual to $\vartheta_D$. Equip $N^\times$ with the flat cylindrical metric $g$ determined by the parallel coframe $\vartheta,\Theta$, and write $\nabla$ for its connection. The end $z\to+\infty$ approaches the zero section.

\begin{lemma}\label{lem:anticanonical-cylindrical-identification}
There is a smooth tubular identification $\Upsilon$ from a neighborhood of the zero section in $N$ to a neighborhood of $D$ in $S$, fixing $D$, such that each fiber disc is holomorphic for $\Upsilon^*J_S$ and
\[
 |\nabla^k(\Upsilon^*\Omega_S-\Omega_0)|_g
 =O(\mathrm e^{-z/2})\qquad(k\geq0).
\]
\end{lemma}
\begin{proof}
Start with a normal exponential map for a smooth Hermitian metric near $D$. Its differential is complex linear along $D$ and induces the identity on the normal quotient. In a parallel local frame and a coordinate $z_0$ with $\vartheta_D=\mathrm dz_0$, choose holomorphic surface coordinates $Z,W$ with $D=\{W=0\}$. The pullbacks under this first map satisfy
\[
 Z=z_0+b(z_0,\bar z_0)w+O(|w|^2),\qquad
 W=w+O(|w|^2).
\]
The residue normalization writes $\Omega_S=(1+Wa(Z,W))\mathrm dW/W\wedge\mathrm dZ$ with $a$ holomorphic. Substitution gives $\Omega_S-\Omega_0=O(|w|)$ in the logarithmic coframe. The possible term containing $\Theta\wedge\mathrm dw$ vanishes. Every derivative by $\partial_{z_0},\partial_{\bar z_0},w\partial_w,\bar w\partial_{\bar w}$ preserves this bound, proving the stated cylindrical estimate for the first map.

It remains to make the fiber discs holomorphic without changing its first jet. On a fiber disc $\Delta_{p,\rho}\subset N_p$, vary the map by $\Phi_Y(q)=\exp_q(Y_q)$. Identify its target tangent spaces with a fixed complex vector space over $p$ by Chern parallel transport along the radial paths and the short exponential paths. In these identifications its antiholomorphic differential is
\[
 \mathcal F_{p,\rho}(Y)
 =\tfrac12\bigl(\mathrm d\Phi_Y+J_S\mathrm d\Phi_YJ_N\bigr)
       |_{T\Delta_{p,\rho}}.
\]
For a fixed Hölder exponent $0<\nu<1$, use rescaled $C^{k,\nu}$ norms on the discs. The Cauchy--Green operator $T_\rho$ gives the normalized right inverse
\[
 C_\rho a=T_\rho a-T_\rho a(0)
                  -w\partial_wT_\rho a(0).
\]
Restrict its domain to $(0,1)$-forms vanishing at the center. Its image consists of fields with zero value and first jet there; on this image $\bar\partial$ and $C_\rho$ are inverse. The differential of $\mathcal F_{p,\rho}$ at zero has the form $\bar\partial+K_{p,\rho}$, with
\[
 \|K_{p,\rho}C_\rho\|\leq C\rho,\qquad
 \|\mathcal F_{p,\rho}(0)\|\leq C\rho^2.
\]
Indeed, the complex structures agree along the zero section, while pulling a one-form back by $w\mapsto\rho w$ supplies the additional factor $\rho$. The inverse is $C_\rho(I+K_{p,\rho}C_\rho)^{-1}$ for small $\rho$. The quantitative implicit-function theorem therefore gives a correction with zero first jet and $Y=O(|w|^2)$.

Compactness of $D$ makes the choice uniform in $p$. The normalized Cauchy operator is equivariant under the unitary changes of fiber frame; the intrinsic transports and uniqueness of the small solution make the corrections agree on overlaps. The parametric implicit-function theorem gives smooth dependence on $p$, and interior elliptic regularity on slightly smaller fiber discs gives smoothness in the fiber. The fixed first jet then gives the quadratic ordinary error, with its cylindrical derivative bounds. Composing the first map with this correction preserves the estimate and proves the lemma.
\end{proof}

\begin{theorem}\label{thm:flat-anticanonical-neighborhood}
\index{Anticanonical bundle!flat elliptic neighborhood}\index{Divisor!flat neighborhood}
Let $S$ be a complex surface and $D\subset S$ a smooth compact connected anticanonical divisor with $D^2=0$. The line bundle $\mathcal O_S(D)$ is unitary flat on a neighborhood of $D$.
\end{theorem}
\begin{proof}
Use the cylindrical identification of \cref{lem:anticanonical-cylindrical-identification}, and omit its pullback notation. A form is horizontal if contraction with $\xi$ and $\bar\xi$ vanishes. For a horizontal form define
\[
 \mathrm d_H\eta=\mathrm d\eta
       -\Theta\wedge\mathcal L_\xi\eta
       -\bar\Theta\wedge\mathcal L_{\bar\xi}\eta.
\]
The degree-$m$ horizontal lift of a form $\eta_m$ valued in $N^{-m}$ is $w^m\eta_m$ in a parallel frame. Direct differentiation identifies $\mathrm d_H$ with its flat covariant derivative. Since the fibers are holomorphic,
\[
 \iota_{\bar\xi}\Omega_S=0,\qquad
 \Omega_S=\Theta\wedge\alpha+\beta,
 \qquad\alpha=\iota_\xi\Omega_S,
\]
with $\alpha,\beta$ horizontal. Closedness gives $\mathcal L_{\bar\xi}\alpha=\mathcal L_{\bar\xi}\beta=0$ and $\mathcal L_\xi\beta=\mathrm d_H\alpha$.

Fix $0<\delta<1/2$, and let $U_T=\{z\geq T\}$. For tensors define the weighted norm
\[
 \|\eta\|_{k,\delta,T}^2
 =\sum_{j=0}^k\int_{U_T}\mathrm e^{2\delta(z-T)}
                    |\nabla^j\eta|_g^2\,\mathrm dV_g.
\]
Let $\mathcal H^{k+1}_{\delta,T}$ be the functions of finite such norm with $\mathcal L_{\bar\xi}f=0$, and let $\mathcal D^k_{\delta,T}$ be the complex two-forms of finite norm with $\mathrm d\gamma=0$ and $\iota_{\bar\xi}\gamma=0$. We construct an operator
\[
 Q_T:\mathcal D^k_{\delta,T}\longrightarrow\mathcal H^{k+1}_{\delta,T},
 \qquad \|Q_T\gamma\|_{k+1,\delta,T}\leq C_{k,\delta}\|\gamma\|_{k,\delta,T},
\]
with constant independent of $T$, satisfying
\begin{equation}\label{eq:anticanonical-horizontal-inverse}
 \mathrm d_H(Q_T\gamma\,\vartheta)
 =\gamma-\Theta\wedge\iota_\xi\gamma,
 \qquad Q_T(\mathrm df\wedge\vartheta)=f.
\end{equation}

Write $\gamma=\Theta\wedge\alpha+\beta$. Its defining equations give $\mathcal L_{\bar\xi}\gamma=0$ and, in the fiberwise homogeneous expansion, $m\beta_m=\mathrm d_H\alpha_m$. Only $m\geq1$ occur: the finite positive weighted norm excludes the zero and negative Laurent modes. Decompose
\[
 \alpha_m=-u_m\vartheta+v_m,
 \qquad u_m\in\Omega^0(D,N^{-m}),\quad
 v_m\in\Omega^{0,1}(D,N^{-m}).
\]
Apply \cref{lem:flat-elliptic-derivative-inverse} to get $w_m$ with $\bar\partial w_m=\iota_v\partial v_m$ and $\|w_m\|_{H^k}\leq\|v_m\|_{H^k}$. Put $f_m=-(u_m+w_m)/m$. Then
\[
 \mathrm d_H(f_m\vartheta)=\beta_m,\qquad
 \|mf_m\|_{H^k}\leq C\|\alpha_m\|_{H^k}.
\]
Also $\|\bar\partial f_m\|_{H^k}\leq C\|\beta_m\|_{H^k}$. Flatness gives the same bound for $\partial f_m$ by integration by parts, and hence controls the pure horizontal derivatives of order $k+1$. A vertical derivative of a homogeneous mode multiplies it by $m$, so every derivative involving a vertical direction is controlled by the preceding bound for $mf_m$. Circle orthogonality and integration of the weight give the required estimate for $f=\sum_{m\geq1}f_m$: the squared radial factor is $\mathrm e^{-mz}$ and $m-2\delta>0$.

For Sobolev coefficients, regularize each mode by the flat base heat operator, which commutes with the defining differential equations. The uniform estimates and orthogonality make the finite sums Cauchy in the target space. Passing to their limit proves the first identity in \eqref{eq:anticanonical-horizontal-inverse}. If $\gamma=\mathrm df\wedge\vartheta$, then $\alpha_m=mf_m\vartheta$, so its $v_m$ and normalized $w_m$ vanish and the construction returns $f_m$. This proves the second identity.

Put $\eta=\Omega_S-\Omega_0$. Choose an integer $k\geq4$. The cylindrical estimate bounds its weighted $H^{k+2}$ norm on $U_{T-1}$ and its first $k+2$ uniform derivatives by a number $a_T\to0$. For small $f\in\mathcal H^{k+1}_{\delta,T}$, the map $\Phi_f(\zeta)=\mathrm e^{f(\zeta)}\zeta$ is a $C^2$ diffeomorphism onto its image and sends $U_T$ into $U_{T-1}$. Its restriction to every fiber is holomorphic. Thus $\Phi_f^*\Omega_S-\Omega_0$ belongs to the actual domain of $Q_T$. Since
\[
 \Phi_f^*\Omega_0-\Omega_0=\mathrm df\wedge\vartheta,
\]
the desired equation becomes
\[
 Q_T(\Phi_f^*\Omega_S-\Omega_0)=f+E_T(f)=0,
 \qquad E_T(f)=Q_T(\Phi_f^*\eta).
\]
On a fixed sufficiently small ball, weighted Sobolev multiplication and composition on the flat cylinder give
\[
 \|E_T(f)\|_{k+1,\delta,T}\leq Ca_T,\qquad
 \|E_T(f)-E_T(g)\|_{k+1,\delta,T}
 \leq Ca_T\|f-g\|_{k+1,\delta,T}.
\]
The constants are independent of $T$: translation in the cylindrical direction is an isometry, the weight is normalized at $T$, and the shifts of $z$ are uniformly bounded. Differentiating the pullback along the segment between $f$ and $g$ takes one Lie derivative of the fixed form $\eta$ and at most one derivative of $f-g$; the gain of one derivative in $Q_T$ restores the indicated norm. For $T$ sufficiently large, $-E_T$ is a contraction of the closed ball of radius $2Ca_T$ into itself. Let $f$ be its fixed point.

The first identity in \eqref{eq:anticanonical-horizontal-inverse} makes the horizontal component of $\Phi_f^*\Omega_S$ vanish. Consequently $\Theta\wedge\Phi_f^*\Omega_S=0$. The pulled-back volume form is nowhere zero and of type $(2,0)$, so $\Theta$ is of type $(1,0)$ for $\Phi_f^*J_S$. Since $\mathrm dz=-(\Theta+\bar\Theta)$ and $\mathrm d\Theta=0$, the function $z$ is pluriharmonic for that complex structure.

Return to $S$ through $A=\Upsilon\circ\Phi_f$. On each fiber, boundedness extends $f$ holomorphically across zero. Its positive-mode expansion gives $f(0)=0$. The uniform small bound on $f$ lets Rouch\'e's theorem compare $w\mathrm e^{f(w)}-a$ with $w-a$ on a fixed smaller circle, for every sufficiently small $a$. Thus the image of $A$ contains a smaller punctured neighborhood of $D$. The function $r=z\circ A^{-1}$ is $C^2$ and pluriharmonic there, hence smooth by elliptic regularity. If $w_S$ is a holomorphic defining function for $D$, the fixed first normal jet of $\Upsilon$ and the bounded factor $\mathrm e^f$ give
\[
 C^{-1}|w|\leq|w_S\circ A|\leq C|w|.
\]
Thus $r+\log|w_S|^2$ is bounded and pluriharmonic off $D$. Both signs have plurisubharmonic extensions across the divisor; their sum is zero almost everywhere and hence everywhere. The extension is therefore pluriharmonic. It follows that $F=-r$, with its logarithmic pole on $D$, satisfies $\ddc F=[D]$ near $D$. Locally, the norm assigning $\exp(F)$ to the squared norm of the defining section is a smooth flat metric, since $F-\log|w_S|^2$ is smooth pluriharmonic. These norms agree on overlaps and give the asserted unitary flat metric.
\end{proof}

\begin{proposition}\label{prop:flat-divisor-rank-one-curvature}
Let $D$ be a smooth nonempty compact divisor in a complex surface $S$, and suppose $\mathcal O_S(D)$ is unitary flat near $D$. It has a smooth semipositive metric on $S$ whose curvature $B$ has rank at most one and vanishes near $D$ and outside a relatively compact neighborhood of $D$. If $S$ is compact K\"ahler, then $\max_S\operatorname{rank}B=1$.
\end{proposition}
\begin{proof}
Let $s_D$ be the defining section and $h_0$ a flat metric on a neighborhood $U$ of $D$. Put $F=\log|s_D|_{h_0}^2$. Choose a relatively compact tubular neighborhood $U_0\Subset U$. Since $F$ tends uniformly to $-\infty$ along $D$ and is finite on $\partial U_0$, adding a constant makes $F>0$ near $\partial U_0$. Choose a smooth concave function $\chi$ with $\chi(t)=t+c$ for $t\leq-1$ and $\chi(t)=0$ for $t\geq0$. Define $\widetilde F=\chi(F)$ on $U_0$ and extend it by zero away from $U_0$. On $S\setminus D$ prescribe
\[
 |s_D|_h^2=\mathrm e^{\widetilde F}.
\]
Near $D$ this gives $h=\mathrm e^c h_0$, so the frame metric extends smoothly and positively across $D$. It also matches smoothly where the cutoff is constant. Poincar\'e--Lelong gives $B=[D]-\ddc\widetilde F$, and off $D$ one has
\[
 B=-\chi''(F)\,\mathrm dF\wedge\dc F\geq0.
\]
This smooth closed form has rank at most one and the asserted vanishing regions. On a compact K\"ahler surface with K\"ahler form $\kappa$, its mixed integral is
\[
 \int_SB\wedge\kappa=\int_D\kappa>0.
\]
Thus it is not identically zero and its maximum rank is one.
\end{proof}

\Needspace{15\baselineskip}
\begin{corollary}\label{cor:flat-divisor-variational-order}
\index{Nef degeneration!flat divisor curvature}
Let $S$ be a compact K\"ahler surface and $D$ a smooth nonempty divisor whose line bundle is unitary flat near $D$. Let $B$ be the curvature from \cref{prop:flat-divisor-rank-one-curvature}. For every fixed smooth $f>0$, every Hermitian form $h$ and every smooth positive probability $\mu$, there are $L,U>0$ such that
\[
 L\epsilon\leq\mathcal V_\mu(fB+\epsilon h)
              \leq U\epsilon,\qquad0<\epsilon\leq1.
\]
For a K\"ahler perturbation $\kappa$, the closed background itself satisfies
\[
 \mathcal V_\mu(B+\epsilon\kappa)
 =2\epsilon\int_D\kappa+\epsilon^2\int_S\kappa^2.
\]
There are choices of $f$ for which $fB$ is nonclosed.
\end{corollary}
\begin{proof}
The inequalities $(\min f)B\leq fB\leq(\max f)B$ and \cref{prop:nef-maximum-rank-comparison} give the two-sided bounds. The polynomial follows from the closed-volume formula and $B^2=0$. At a point where $B$ has rank one, choose a smooth positive $f$ whose differential has a nonzero component tangential to the kernel of $B$. Then $\mathrm d(fB)=\mathrm df\wedge B\ne0$ there.
\end{proof}

\Needspace{12\baselineskip}
\begin{proposition}\label{prop:nine-point-anticanonical-curvature}
\index{Anticanonical bundle!nine-point blowup}
Let $S$ be obtained from $\mathbb P^2$ by nine successive point blowups, each centered on the strict transform of a fixed smooth cubic. Its final strict transform $D$ is a smooth elliptic anticanonical curve with $D^2=0$. The anticanonical bundle has a smooth semipositive metric with rank-one curvature as above. It is semiample exactly when $N=\mathcal O_D(D)$ is torsion. If $N$ has exact order $m$, then $|mD|$ is basepoint-free and no smaller positive multiple moves. If $N$ is non-torsion, then $|kD|=\{kD\}$ for every $k\geq1$.
\end{proposition}
\begin{proof}
A blowup at a smooth point of the cubic replaces its class by its pullback minus the exceptional curve, exactly as for the anticanonical class. The strict transform remains isomorphic to the cubic, and its self-intersection decreases by one. After nine steps $D^2=9-9=0$. The flat neighborhood theorem and curvature construction apply.

For $k\geq1$, use the restriction sequences
\[
 0\longrightarrow\mathcal O_S((k-1)D)
 \longrightarrow\mathcal O_S(kD)
 \longrightarrow N^k\longrightarrow0.
\]
A nontrivial degree-zero line on an elliptic curve has no section: any nonzero section would have a zero divisor of degree zero and trivialize the line. If $N$ is non-torsion, induction gives $h^0(S,\mathcal O_S(kD))=h^0(S,\mathcal O_S)=1$, proving the singleton systems and failure of semiampleness.

If $N$ has exact order $m$, then $H^0(D,N^k)=H^1(D,N^k)=0$ for $1\leq k<m$, the latter by Serre duality on the elliptic curve. Since $S$ is rational, $H^1(S,\mathcal O_S)=0$. The same sequences give $h^0(S,kD)=1$ and $H^1(S,\mathcal O_S(kD))=0$ for $0\leq k<m$. The constant section of $N^m\simeq\mathcal O_D$ therefore lifts to a section of $\mathcal O_S(mD)$ nonzero everywhere on $D$. Together with $s_D^m$, it generates that line globally. Thus $|mD|$ is basepoint-free, and no smaller multiple moves.

Non-torsion configurations exist with nine distinct centers. Fix the first eight points on the cubic and vary the ninth. Under the unchanged identification with that elliptic curve, the normal line is $\mathcal O_D(3)\otimes\mathcal O_D(-p_1-\cdots-p_9)$. Varying $p_9$ parametrizes its degree-zero Picard group by the elliptic Abel map. The torsion subgroup is countable, while the curve is uncountable; excluding the first eight points still leaves non-torsion choices. Hence \cref{cor:flat-divisor-variational-order} supplies nonclosed backgrounds of order one even when the associated anticanonical line is not semiample.
\end{proof}
